\section{All geometric signs when degree divided by dimension is bounded}\label{sec:small-cap-bounded-degree-ratio}

Let $d\to\infty$, $b=d^2/n\to\infty$, $b=o(d^{1/3})$, and $1\le r\le d$ with $r=O(b)$. For the equal-column capped-row array counts $a_i$ and degree $D=\lfloor nr/d\rfloor$, every internal geometric point $t_k=\sqrt{a_{k-1}/a_{k+1}}$, $1\le k<D$, eventually satisfies $(-1)^kP(-t_k)>0$, uniformly along these sequences. In particular all roots are simple and negative. The degree-zero and degree-one cases are immediate.

Use the \href{https://github.com/ProgramComputer/open-math-research/blob/4fb624ba937a8eecc1361a6066185cd81ab7b172/topics/balanced-arrays/orientation-normalization.md}{orientation bound}
$\Gamma_i=a_i^2/(a_{i-1}a_{i+1})\ge(1+1/i)^d(1+1/(n-i))^d$. Thus $t_i$ are strictly increasing, $\Gamma_i>4$ suffices for the corresponding sign by pairing both decreasing alternating tails after their first term, and $k/d\to0$ is covered. Also $D/d\le r/b=O(1)$. It remains to examine subsequences with $k/d\to\beta\in(0,\infty)$ and any fixed feasible shifts of these quotas.

\subsection{Uniform margin correction}

Let $H(z)=\sum_{s=0}^r\binom ds z^s$, $U_i=[z^{di}]H(z)^n$, $K_i=\binom ni^d/\binom{nd}{di}$, $\mu=di/n\asymp b$, and $p=i/n=\mu/d=o(1)$. At a feasible integer quota $i=k+O(1)$ let $X$ have the binomial size law with power tilt and cap $r$ tuned to mean $\mu$, and set $V=\mathrm{Var}(X)$. At exact saturation take $X=r$. Then uniformly on these shifts,

\begin{equation}
a_i=K_i U_i\exp\!\left(-\frac12+\frac{V}{2\mu(1-p)}+o(1)\right).
\tag{L1}
\end{equation}

Indeed for any prescribed row profile $(s_j)$ the total $S=di\asymp d^2$, maximum product is at most $ri=O(bd)=o(S^{2/3})$, and with column margins all $i$ the scaled column variance is zero. The scaled row variance is $\sum_j(s_j-\mu)^2/[n\mu(1-p)]\le(r-\mu)/(1-p)=O(b)$, satisfying the additional $O(S^{1/3})$ requirement. Greenhill--McKay--Wang Corollary 5.2 gives the profile count as $K_i\prod_j\binom d{s_j}$ times $\exp[-1/2+\sum_j(s_j-\mu)^2/(2n\mu(1-p))+O(b/d^{1/3})]$, uniformly. \href{https://web.maths.unsw.edu.au/~csg/papers/irregular-final.pdf#page=33}{GMW final text, Corollary 5.2}, \href{https://doi.org/10.1016/j.jcta.2005.03.005}{published DOI}.

Conditional on $\sum_j X_j=di$, independent tuned rows have precisely the product binomial weights after normalization by $U_i$. Writing $Y_j=(X_j-\mu)^2/[\mu(1-p)]$, one has $0\le Y_j\le Cb$ and $\mathbb EY_j\le1$ by the success/failure bias variance bound $V\le\mu(1-\mu/d)$. The log-concavity of the row law and its convolutions gives $\mathbb P(\sum_jX_j=di)\ge1/(nr+1)$ by Jensen for the concave interpolation of the log lattice mass at its mean, followed by the entropy bound. The chord bound for $e^x$ gives $\mathbb E e^{Y_j/n}\le1+C/n$ since $b/n=o(1)$, hence $\mathbb E e^{\overline Y}=O(1)$. With $\varepsilon=d^{-1/12}$, Hoeffding gives $\mathbb P(|\overline Y-\mathbb E Y_j|>\varepsilon)\le2e^{-c n\varepsilon^2/b^2}$, superpolynomial in $d$ since $n=d^2/b$ and $b^3=o(d)$. Cauchy--Schwarz and the polynomial conditioning cost show the exceptional contribution to $\mathbb E(e^{\overline Y/2}\mid\sum_jX_j=di)$ is $o(1)$. On the typical event it is $e^{\mathbb E Y_j/2}(1+o(1))$. The deterministic saturation case is immediate. This proves (L1). The elementary binomial prefactor curvature satisfies $|2\log K_i-\log K_{i-1}-\log K_{i+1}|=O(d/i^2+d/(n-i)^2)=o(1)$: compare each of the $d$ denominator factors in adjacent binomial ratios with the limiting factors at quota $i$ or $n-i$, as in \href{https://github.com/ProgramComputer/open-math-research/tree/ba7ddfd7efb37cac33769f4fdfed749e1b5bf58a/topics/balanced-arrays/dense-regime/rank-dependent}{the rank-dependent result}.

\subsection{Curvatures and signs}

For real quota $u$ tune the size law at mean $\mu(u)=du/n$ and set $h=r-\mu$. Let $E(u)=n\log H(z)-du\log z$. At integer quotas $U_i=e^{E(i)}L_i$, where $L_i$ is the integer-mean probability for the sum of independent tuned rows. In the interior $E''(u)=-b/V(u)$. The log-concave lattice estimates \hyperref[sec:subquadratic-scaling-corank-upper]{at positive variance} give $\left|dV/d\mu\right|\le C(1+\sqrt V)$ and, if $n\min(1,V)\to\infty$, $L_i=(2\pi nV(i))^{-1/2}(1+o(1))$. Thus $\sqrt{1+V}$ is Lipschitz in the mean, which moves by $b/d=o(1)$ at unit quota shifts.

The deficit $J=r-X$ has a log-concave mass. Dividing its masses by those of a geometric law of the same mean $h$ gives a log-concave ratio, so the positive part of the probability-mass difference is supported on an interval. Both total masses and means agree. Subtracting the affine chord of a convex test function across the crossing boundaries proves convex-order domination (if the laws differ the equal means force a crossing on both sides). In particular $V\le h(1+h)$. Also $V\ge c\min(1,\mu,h)$ by the \href{https://github.com/ProgramComputer/open-math-research/blob/4fb624ba937a8eecc1361a6066185cd81ab7b172/topics/balanced-arrays/quadratic/very-tall-quadratic-interior-saddle.md}{modal-atom argument}.

Pass to a subsequence with $V(k)/b\to v\in[0,\infty)$, available since $V(k)\le\mu(k)=O(b)$.

\begin{itemize}
\item If $v>0$, then $h\gtrsim\sqrt b$ by the geometric variance bound. The lower support distance is $k\asymp d$, and the upper support distance is at least $c d/\sqrt b\to\infty$. At any fixed shifts, $V(u)/b\to v$ uniformly by the derivative bound. The second log difference of $E$ therefore tends to $1/v$ in positive-curvature convention, and the lattice prefactors have log differences $o(1)$. The correction in (L1) likewise varies by $o(1)$ across fixed shifts, since $V$ changes by at most $O((1+\sqrt b)b/d)=o(b)$ and $\mu\asymp b$. Thus $\log\Gamma_{k+j}\to1/v$ for every fixed integer $j$.
\item If $v=0$ and $h\ge\eta>0$ along a subsequence, then $V\ge c_\eta>0$ at fixed shifts, because $\mu\asymp b$ and $b/d=o(1)$. The variance changes by $o(1)$ in absolute value: $(1+\sqrt b)b/d=o(1)$. The lattice prefactors have log differences $o(1)$, and the curvature of $E$ tends to infinity since $V=o(b)$ uniformly across the neighboring real quotas. The correction in (L1) tends to $-1/2$. Hence $\Gamma_k\to\infty$.
\item If $h\to0$, use the \href{https://github.com/ProgramComputer/open-math-research/blob/04564a728c21eab6ce511b98a48ccaac754648ff/topics/balanced-arrays/very-tall/very-tall-upper-quota.md#the-last-od-quotas}{scalar negative-binomial step inequality} with deficit $e=nh\ge d$ and row deficit weights $\binom d{r-j}$. It yields
\end{itemize}
  $\log(U_k^2/(U_{k-1}U_{k+1}))\ge[d(n-1)/(n+e+d-1)]\log(1+d/e)\ge(1-o(1))b/(h+b/d)\to\infty$.
  This includes the exact saturation neighbor. Here $V=o(1)$ at all three quotas by $V\le h(1+h)$ and $b/d=o(1)$, so (L1) and the prefactor give $\Gamma_k\to\infty$.

The divergent-curvature cases give the sign immediately from $\Gamma_k>4$. In the finite-curvature case the absolute normalized term at each fixed offset $j$ at $-t_k$ tends to $e^{-j^2/(2v)}$, by multiplying the fixed-offset ratios. The orientation bound gives uniform curvature at least $d/(D+1)\ge c>0$ at every internal index, so all normalized terms are bounded by $e^{-c j^2/2}$, setting missing terms to zero. Dominated convergence gives the alternating Gaussian sum, strictly positive by Poisson summation. This proves the sign along every subsequence, hence uniformly as claimed. Alternating signs at zero, at the ordered internal points, and at large negative argument give the full degree of distinct simple negative roots.
